\section{地磁気偏角}

\subsection{令和元年度 問27}
\label{q:r01-reg-applied-27}

\noindent\textbf{原典：} sources/R01/17\_applied\_science.pdf、PDF物理9頁。

\textbf{公式正答：} 2（single）。

\paragraph{問題の要旨} 令和元年度試験当時の日本周辺の偏角の方向と地域差から適切な記述を選ぶ。

\paragraph{問われる核心} 地磁気偏角の方向・分布・時点

\paragraph{必要知識} 真北と磁北、偏角の符号、二次元地理分布、地磁気永年変化、磁気図の元期。

\paragraph{考え方（査読待ち初稿）} 偏角は水平磁場が真北となす角である。原典の「現在」を令和元年度時点とし、それ以前に公表された2015.0年磁気図を確認する。日本周辺は西偏で、沖縄本島の等偏角線は約5度である。緯度・経度の片方だけで同値になる量ではないので第2肢を選ぶ。

\paragraph{正答を選ぶ理由} 偏角は水平磁場が真北となす角である。原典の「現在」を令和元年度時点とし、それ以前に公表された2015.0年磁気図を確認する。日本周辺は西偏で、沖縄本島の等偏角線は約5度である。緯度・経度の片方だけで同値になる量ではないので第2肢を選ぶ。

\paragraph{選択肢の検討（査読待ち初稿）}\mbox{}\par

\noindent\textbf{1：不適切} 全般的に北ほど偏角が小さい。当時の2015.0年磁気図では沖縄本島付近約5度西偏に対し、北海道では概ね9〜10度西偏であり、北ほど小さいという傾向と合わない。\par

\noindent\textbf{2：適切} 沖縄本島の偏角は約5度。当時参照可能な国土地理院の2015.0年磁気図で、沖縄本島付近は約5度の西偏を示す。概略値として適切である。\par

\noindent\textbf{3：不適切} 真北から東へ偏っている。日本周辺は西偏であり、東西が逆である。偏角の符号規約と方位を区別する。\par

\noindent\textbf{4：不適切} 同じ経度なら偏角は同じ。磁場は経度だけでは決まらず、同じ経度でも緯度に伴って偏角が異なる。\par

\noindent\textbf{5：不適切} 同じ緯度なら偏角は同じ。非軸対称な磁場があるため同緯度でも経度によって変わる。\par

\paragraph{再利用する考え方} 地磁気偏角の方向・分布・時点を、測定量・推定対象・基準面・近似条件に分けて理解する。名称だけで探査法を対応させない。

\paragraph{共通理論} \hyperref[theory:geomagnetism:declination]{第\ref*{theory:geomagnetism:declination}節}

\par 参考：\cite{gsi_declination2015}。

\par 難易度 3（校準初稿）、優先度未評価。
